\chapter*{Executive Summary}
\addcontentsline{toc}{chapter}{Executive Summary}

This project investigates whether multi-way representations of resting-state
\EEG{} can support three-class subject-level discrimination among Alzheimer's
disease (\AD), frontotemporal dementia (\FTD), and cognitively normal controls
(\CN). The source is the public OpenNeuro \code{ds004504} cohort: 88 people,
19 scalp channels, 500 Hz eyes-closed recordings, and author-supplied artifact-
cleaned derivatives \citep{miltiadous2023dataset,openneuro2026dataset}.

The work evolved through four experiments. It began with a deliberately balanced
69-person primary cohort and exploratory non-negative spectral tensor methods.
It then corrected channel labels and EEGLAB boundary handling, introduced
training-only TT/Tucker bases and nested subject cross-validation, expanded to all
88 people using raw-waveform delay moments plus class/subject weighting, and
finally performed a bounded 139-configuration comparison of logistic regression,
RBF SVM, random forest, XGBoost, and supervised TT variants.

\begin{keybox}{Headline outcome}
The strongest observed classifier-family result is the tuned RBF SVM at
\textbf{59.7\% mean balanced accuracy} (repeat SD 0.9 percentage points), with
57.4\% \AD{} recall, 43.5\% \FTD{} recall, and 78.2\% \CN{} recall. The V4
predeclared complete-pipeline endpoint is \textbf{54.5\%}. This is useful course
evidence above the 33.3\% single-class balanced-accuracy baseline, but it is not
an independently confirmed 60\% clinical result.
\end{keybox}

The central empirical lesson is that tensorization alone did not guarantee a
gain. Corrected V2 Tucker/HOSVD reached 56.0\% as the strongest method family,
while V3's same-cohort spectral control (55.8\%) exceeded every waveform-only
branch (49.9--51.8\%). V4's TT/spectral fusion family reached 59.1\%, but
training-only feature selection could favor spectral predictors, so the result
does not isolate a causal TT contribution. Across versions, \FTD{} remains the
hardest class and the principal data limitation is 23 independent \FTD{} people,
not merely the number of four-second windows.

The pipeline is methodologically stronger than its first prototype because all
current learned stages are fitted on training people only; windows from a person
never cross validation boundaries; per-person pooling prevents long recordings
from becoming pseudo-replicates; class weights equalize diagnosis-level loss;
and verification replays saved predictions, checks split membership, confirms
training-only bases/scalers, and fingerprints source arrays.

Figure~\ref{fig:executive-pipeline} summarizes how these stages connect.
\begin{figure}[!htbp]
\centering
\begin{tikzpicture}[node distance=8mm and 7mm,
every node/.style={font=\small},
block/.append style={text width=43mm,minimum height=15mm},
process/.append style={text width=43mm,minimum height=15mm}]
\node[block] (source) {OpenNeuro\\88 people, 19 channels};
\node[process,right=of source] (qc) {Duration audit\\CAR, QC, boundary guards};
\node[block,right=of qc] (rep) {Spectral tensors or\\waveform delay moments};
\node[process,below=of rep] (cv) {Nested participant CV\\training-only basis and model};
\node[block,left=of cv] (out) {AD / FTD / CN\\one prediction per person};
\draw[arrow] (source)--(qc);
\draw[arrow] (qc)--(rep);
\draw[arrow] (rep)--(cv);
\draw[arrow] (cv)--(out);
\node[align=left,text width=43mm,below=of source] {Each person's windows stay together throughout validation.};
\end{tikzpicture}
\caption{Completed NeuroTensor workflow. Predictive evaluation is performed at the participant level.}
\label{fig:executive-pipeline}
\end{figure}


\clearpage
